\documentclass[11pt]{article}
\usepackage[T1]{fontenc}
\usepackage[utf8]{inputenc}
\usepackage{lmodern,amsmath,amssymb,amsthm,microtype}
\usepackage[a4paper,margin=28mm]{geometry}
\usepackage[colorlinks=true,linkcolor=blue,citecolor=blue,urlcolor=blue]{hyperref}
\hypersetup{
  pdftitle={A sqrt(3) growth criterion for weighted reciprocal series},
  pdfauthor={twinheterodox},
  pdfsubject={Irrationality under irregular growth for reciprocal products with weights (2,1)},
  pdfcreator={LaTeX}
}
\newtheorem{theorem}{Theorem}[section]
\newtheorem{lemma}[theorem]{Lemma}
\numberwithin{equation}{section}
\newcommand{\Q}{\mathbb{Q}}
\newcommand{\N}{\mathbb{N}}
\title{A $\sqrt{3}$ growth criterion for weighted reciprocal series}
\author{twinheterodox\\[2pt]\normalsize Prepared with Codex}
\date{13 November 2026}

\begin{document}
\maketitle

\begin{abstract}
We prove that $\sum_{n\ge1}1/(a_n^2a_{n+1})$ is irrational for every
strictly increasing sequence of positive integers satisfying
\[
  \limsup_{n\to\infty}\frac{\log a_n}{(\sqrt3)^n}=\infty.
\]
This lowers the sufficient growth base from $2$ to $\sqrt3$ in the
criterion of Barreto, Kang, Kim, Kova\v{c} and Zhang
\cite[Theorem~3]{BKKKZ} for weights $(2,1)$.
Irregular growth is allowed. The result also extends to nondecreasing
sequences and positive integer numerators subject to polynomial size
bounds.
\end{abstract}

\section{Introduction}

Growth criteria for irrationality ask how fast an integer sequence must
grow before an associated reciprocal series is forced to have an
irrational sum. Such a criterion applies to entire families of sums,
without a closed form for the sum or information about the prime factors
of the sequence. The underlying constraint is arithmetic: a rational
partial sum cannot approach a different rational number more closely
than their denominators permit.

Barreto, Kang, Kim, Kova\v{c} and Zhang \cite{BKKKZ} answered a question
of Erd\H{o}s and Graham \cite[p.~64]{ErdosGraham} about reciprocal
products of successive sequence terms and developed criteria for
more general products
$a_n^{w_0}\cdots a_{n+d-1}^{w_{d-1}}$. When some weights exceed one,
their sufficient growth condition and their rational examples can
leave a gap \cite[Section~5]{BKKKZ}. We improve the sufficient
condition for the two adjacent factors with weights $(2,1)$.

\begin{theorem}\label{thm:strict}
Let $(a_n)_{n\ge1}$ be a strictly increasing sequence of positive
integers. If
\begin{equation}\label{eq:growth}
  \limsup_{n\to\infty}\frac{\log a_n}{(\sqrt3)^n}=\infty,
\end{equation}
then
\[
  \sum_{n=1}^{\infty}\frac{1}{a_n^2a_{n+1}}
\]
is irrational.
\end{theorem}

All logarithms are natural. Condition~\eqref{eq:growth} is equivalent
to $\limsup_{n\to\infty}a_n^{1/(\sqrt3)^n}=\infty$. It requires large
values along a subsequence; it imposes no limit on the ratios of
successive logarithms.

\subsection{Comparison with the earlier bound}

For weights $(2,1)$, the polynomial in \cite[Theorem~3]{BKKKZ} is
\[
  (x-1)(2+x)-2x=x^2-x-2,
\]
whose positive root is $2$. That theorem therefore requires
$\limsup a_n^{1/2^n}=\infty$. Theorem~\ref{thm:strict} replaces $2$ by
$\sqrt3$, and hence applies under a weaker growth condition.
For unit numerators, the base-$2$ condition also follows from
Erd\H{o}s's criterion \cite[Theorem~1]{Erdos1975}, applied to
$x_n=a_n^2a_{n+1}$.

For a concrete comparison, suppose $\log a_n\sim\lambda^n$. The new
criterion applies for every $\lambda>\sqrt3$, whereas the earlier
criterion requires $\lambda>2$ within this family. In particular,
\[
  a_n=\left\lceil\exp\bigl((9/5)^n\bigr)\right\rceil
\]
is covered by Theorem~\ref{thm:strict}, while
$a_n^{1/2^n}\to1$. The same theorem also covers sequences whose growth
is irregular; the proof must allow long stretches of slower growth
between the large values supplied by \eqref{eq:growth}.

The rational examples in \cite[Theorem~5]{BKKKZ} are governed here by
\[
  (x-1)(2+x)-x=x^2-2.
\]
For every $C>1$, that theorem gives a strictly increasing integer
sequence for which $a_n^{1/(\sqrt2)^n}\to C$ and the reciprocal series
has a rational sum. Thus the sufficient base is now $\sqrt3$, while
the rational constructions occur at $\sqrt2$. The optimal threshold
remains open.

\subsection{A general form}

We prove a version that also permits repetitions in $(a_n)$ and
positive integer numerators.

\begin{theorem}\label{thm:general}
Let $(a_n)_{n\ge1}$ be nondecreasing positive integers and let
$(b_n)_{n\ge1}$ be positive integers. Suppose that, for some real
$0\le\eta<\tau$, one has
\begin{equation}\label{eq:size}
  b_n\le n^\eta,\qquad a_n^2a_{n+1}\ge n^{1+\tau}
\end{equation}
for all sufficiently large $n$. If \eqref{eq:growth} holds, then the
convergent series
\begin{equation}\label{eq:series}
  S=\sum_{n=1}^{\infty}\frac{b_n}{a_n^2a_{n+1}}
\end{equation}
is irrational.
\end{theorem}

Theorem~\ref{thm:strict} follows by taking $b_n=1$, $\eta=0$ and
$\tau=2$, since strict increase gives $a_n\ge n$.
The size assumptions in Theorem~\ref{thm:general} are the same type
as those in \cite[Theorem~3]{BKKKZ}: they give a uniform bound for
tails even where growth is slow.

\subsection{The source of \texorpdfstring{$\sqrt3$}{sqrt(3)}} 

The improvement begins with the denominator of a partial sum.
The first three terms of \eqref{eq:series} have denominators
\[
  a_1^2a_2,\qquad a_2^2a_3,\qquad a_3^2a_4.
\]
Their sum is cleared by $a_1^2a_2^2a_3^2a_4$: only the first power
of the last factor is needed. For the first $N$ terms, the corresponding
common denominator is
\begin{equation}\label{eq:clearer}
  D_N=\left(\prod_{i=1}^N a_i^2\right)a_{N+1}.
\end{equation}
Writing $u_n=\log a_n$ and $V_n=\sum_{i=1}^n u_i$, we obtain
\begin{equation}\label{eq:endpoint}
  \log\frac{D_N}{a_{N+1}^2a_{N+2}}=3V_N-V_{N+2}.
\end{equation}
The relation $V_{N+2}=3V_N$ has growth scale $(\sqrt3)^N$.
This explains the base in the theorem.

Identity~\eqref{eq:endpoint} only compares the common denominator with
the first omitted term. For irregular growth, the rest of the tail
need not be negligible. We adapt the tail estimates and the argument
for selecting a peak after a slow index in
\cite[Lemmas~9--11 and Proposition~12]{BKKKZ}. These build on the
tail and peak arguments of Erd\H{o}s \cite[pp.~3--6]{Erdos1975}.
The two changes are the endpoint denominator \eqref{eq:clearer} and
the use of peaks of the cumulative logarithms $V_n$. These changes
preserve the saving in \eqref{eq:endpoint} while controlling the entire
tail. We give all the required estimates below.

\section{Two uniform tail estimates}\label{sec:tails}

The proof uses a coarse estimate valid everywhere and a sharper estimate
on intervals of fast growth. The endpoints of those intervals will be
chosen later, so the constants must be independent of them. The following
estimates adapt \cite[Lemma~10]{BKKKZ}; the corresponding bounds with
unit numerators appear in \cite[pp.~3--4]{Erdos1975}.

\begin{lemma}\label{lem:tails}
Let $(x_n)$ be a nondecreasing sequence of positive real numbers and let
$0\le y_n\le n^\eta$. Fix $0\le\eta<\tau$, $c_0>0$ and a positive
integer $A$.
\begin{enumerate}
\item If $x_n\ge c_0 n^{1+\tau}$ for every $n\ge A$ and $X=x_A\ge1$,
then
\begin{equation}\label{eq:polytail}
  \sum_{n\ge A}\frac{y_n}{x_n}
  \le C_{\mathrm{pol}}X^{-\alpha},
  \qquad \alpha=\frac{\tau-\eta}{1+\tau}>0.
\end{equation}
The constant $C_{\mathrm{pol}}$ depends only on $\eta,\tau,c_0$.
\item If $x_n\ge e^n$ for $A\le n\le B$, where $B\ge A$ is an integer
or $B=\infty$, then
\begin{equation}\label{eq:exptail}
  \sum_{n=A}^{B}\frac{y_n}{x_n}
  \le C_{\mathrm{exp}}\frac{(\log x_A)^{1+\eta}}{x_A}.
\end{equation}
The constant $C_{\mathrm{exp}}$ depends only on $\eta$.
\end{enumerate}
Neither constant depends on the endpoints or on the particular
sequences.
\end{lemma}

\begin{proof}
Put $X=x_A$ and group the relevant indices according to
$2^jX\le x_n<2^{j+1}X$, for $j\ge0$.
In the first case, every index in the $j$th group is at most
$(2^{j+1}X/c_0)^{1/(1+\tau)}$. Since $\eta\ge0$, summing $n^\eta$
up to this bound shows that the group's contribution is at most
\[
  C\,\frac{(2^{j+1}X)^{(1+\eta)/(1+\tau)}}{2^jX}.
\]
The sum over $j$ is geometric with ratio $2^{-\alpha}<1$, proving
\eqref{eq:polytail}.

In the second case, every index in the $j$th group is at most
$\log X+(j+1)\log2$. Its contribution is therefore at most
\[
  C\,\frac{(\log X+j+1)^{1+\eta}}{2^jX}.
\]
Here $\log X\ge A\ge1$. Summing the convergent series
$\sum_{j\ge0}(j+2)^{1+\eta}/2^j$ proves \eqref{eq:exptail}.
The same bounds apply when an interval ends partway through a group.
\end{proof}

\section{Proof of Theorem~\ref{thm:general}}\label{sec:proof}

\subsection{What rationality would force}

Write
\[
  x_n=a_n^2a_{n+1},\qquad
  T_N=\sum_{n>N}\frac{b_n}{x_n}.
\]
The series converges because \eqref{eq:size} bounds its terms by
$n^{-1-(\tau-\eta)}$ for all large $n$.
We use the elementary criterion in \cite[Lemma~8]{BKKKZ}.
Suppose, for a contradiction, that $S=p/q$ with $q$ a positive integer.
The integer $D_N$ in \eqref{eq:clearer} clears each term of the partial
sum through $N$. Thus $qD_NT_N$ is a positive integer, and
\begin{equation}\label{eq:integerlower}
  D_NT_N\ge\frac1q\qquad(N\ge1).
\end{equation}
We will find indices at which this same cleared tail tends to zero.

Retain $u_n=\log a_n$ and $V_n=\sum_{i=1}^n u_i$, with $V_0=0$, and put
\begin{equation}\label{eq:EF}
  E_N=\log D_N=V_N+V_{N+1},\qquad
  F_N=\log x_{N+1}=2u_{N+1}+u_{N+2}.
\end{equation}
The endpoint identity becomes
\begin{equation}\label{eq:EFidentity}
  E_N-F_N=3V_N-V_{N+2}.
\end{equation}

By \eqref{eq:polytail}, with $c_0=1$ and
$\alpha=(\tau-\eta)/(1+\tau)$, the full tail satisfies
$T_N\le C_{\mathrm{pol}}e^{-\alpha F_N}$ for all large $N$.
Combining this with \eqref{eq:integerlower} gives
\[
  \alpha F_N\le E_N+\log(C_{\mathrm{pol}}q).
\]
Monotonicity and \eqref{eq:size} imply $a_n\to\infty$, so $E_N\to\infty$.
Consequently, for all sufficiently large $N$,
\begin{equation}\label{eq:jump}
  F_N\le\frac{2}{\alpha}E_N,\qquad
  E_{N+1}\le H E_N,\qquad H=1+\frac{2}{\alpha}.
\end{equation}
Indeed, $E_{N+1}-E_N=u_{N+1}+u_{N+2}\le F_N$.

The bound \eqref{eq:jump} is deliberately coarse. Under the assumption
of rationality, it prevents a slow index from being followed immediately
by an uncontrolled jump. It also gives $E_N=O(H^N)$ and hence
\begin{equation}\label{eq:logF}
  \log(F_N+1)=O(N+1).
\end{equation}
All constants from this point may depend on the fixed sequences,
$\eta,\tau$ and the hypothetical denominator $q$, but not on the
indices that will be selected.

\subsection{Eventual exponential growth}

Set $c=\sqrt3$ and $W_n=V_n/c^n$. Suppose first that $a_n\ge e^n$
for every sufficiently large $n$. Then $x_n\ge e^n$ there, and
\eqref{eq:exptail} applied to the whole tail gives
\[
  T_N\le C_{\mathrm{exp}}F_N^{1+\eta}e^{-F_N}.
\]
Using \eqref{eq:integerlower}, \eqref{eq:EFidentity} and \eqref{eq:logF},
we obtain
\[
  V_{N+2}\le 3V_N+(1+\eta)\log F_N+O(1)
           \le 3V_N+C(N+1)
\]
for a fixed $C>0$ and all large $N$. Since $c^2=3$,
\[
  W_{N+2}\le W_N+\frac{C(N+1)}{c^{N+2}}.
\]
The errors have a finite sum on each parity class. Iterating separately
on the even and odd indices therefore bounds $W_n$.
But $V_n\ge u_n$, so this contradicts \eqref{eq:growth}.

\subsection{A first peak after a slow index}

It remains to treat sequences with $a_n<e^n$ at infinitely many
indices. We will split a tail inside a fast stretch. A first peak
will control the near part, and a later large value will control the
part beyond the stretch.

Choose increasing indices $R$ such that
\begin{equation}\label{eq:mu}
  \mu_R=\frac{u_R}{c^R}\longrightarrow\infty.
\end{equation}
For large $R$, these indices are fast: $a_R\ge e^R$.
Let $P=P(R)<R$ be the last index before $R$ with $a_P<e^P$.
It exists for all large $R$. For every fixed slow index $t$, the
selected $R$ eventually exceed $t$, and then $P(R)\ge t$. Since
the slow indices are unbounded, $P(R)\to\infty$. By definition,
\begin{equation}\label{eq:faststretch}
  a_n\ge e^n\qquad(P<n\le R).
\end{equation}

Monotonicity gives $V_P\le P u_P<P^2$ and $E_{P-1}\le2P^2$.
For all large $P$, \eqref{eq:jump} therefore yields
\[
  V_{P+1}\le E_P\le2H P^2.
\]
It follows that $W_P<1$ and $W_{P+1}<1$. On the other hand,
$W_R\ge\mu_R\to\infty$, so $R\ge P+2$ eventually.

We need a peak that produces a definite saving in
\eqref{eq:EFidentity}, while keeping the earlier denominator under
control. Let $L=e^2$ and choose the first $Q\in[P+2,R]$ such that
\begin{equation}\label{eq:peak}
  W_Q>(1+Q^{-2})\max(1,W_P,\ldots,W_{Q-1}).
\end{equation}
Such a $Q$ exists for all large selected $R$. Otherwise the successive
running maxima, starting from $1$, could increase only by the factors
$1+j^{-2}$, and would remain at most
\[
  \prod_{j\ge1}(1+j^{-2})
  \le\exp\left(\sum_{j\ge1}j^{-2}\right)\le L.
\]
This contradicts $W_R\ge\mu_R\to\infty$.
The summable-increment argument appears in \cite[p.~4]{Erdos1975};
we follow its use for weighted series in
\cite[Lemma~9 and Proposition~12]{BKKKZ}.

Because $Q$ is the first such peak,
\begin{equation}\label{eq:beforepeak}
  W_j\le L\quad(P\le j<Q),\qquad
  E_{Q-2}=V_{Q-2}+V_{Q-1}\le2Lc^Q.
\end{equation}
Also $V_Q>c^Q$, and comparison with $Q-2$ in \eqref{eq:peak} gives
\begin{equation}\label{eq:saving}
  E_{Q-2}-F_{Q-2}=3V_{Q-2}-V_Q
  \le-\frac{V_Q}{Q^2+1}
  \le-\frac{c^Q}{Q^2+1}.
\end{equation}
Here $Q-2\ge P$, so the earlier index is included in the maximum.
Moreover, $Q\to\infty$ because $Q\ge P+2$.

\subsection{Both parts of the cleared tail tend to zero}

Split the tail after $Q-2$ as
\begin{equation}\label{eq:split}
  D_{Q-2}T_{Q-2}
  =D_{Q-2}\sum_{n=Q-1}^{R}\frac{b_n}{x_n}
   +D_{Q-2}T_R.
\end{equation}
The bound on $E_{Q-2}$ in \eqref{eq:beforepeak} has two uses: it
controls the logarithmic factor in the near-tail estimate, and it
limits the clearing cost for the far tail.

For the near part, \eqref{eq:faststretch} gives $x_n\ge e^n$
throughout $Q-1\le n\le R$. By \eqref{eq:jump} and
\eqref{eq:beforepeak},
\[
  F_{Q-2}\le\frac{2}{\alpha}E_{Q-2}
          \le\frac{4L}{\alpha}c^Q.
\]
The finite-interval bound \eqref{eq:exptail} and
\eqref{eq:saving} now give
\begin{align}
  \log\left(D_{Q-2}\sum_{n=Q-1}^{R}\frac{b_n}{x_n}\right)
  &\le E_{Q-2}-F_{Q-2}+(1+\eta)\log F_{Q-2}+O(1)\notag\\
  &\le-\frac{c^Q}{Q^2+1}+(1+\eta)Q\log c+O(1)
       \longrightarrow-\infty. \label{eq:near}
\end{align}
The constants are independent of the moving endpoints $P,Q,R$.

For the far part, use \eqref{eq:polytail} and
$x_{R+1}=a_{R+1}^2a_{R+2}\ge a_R^3$. Together with $Q\le R$,
\eqref{eq:beforepeak} and \eqref{eq:mu}, this gives
\begin{align}
  D_{Q-2}T_R
  &\le C_{\mathrm{pol}}\exp(E_{Q-2}-\alpha\log x_{R+1})\notag\\
  &\le C_{\mathrm{pol}}\exp(2Lc^Q-3\alpha u_R)\notag\\
  &\le C_{\mathrm{pol}}\exp\bigl((2L-3\alpha\mu_R)c^R\bigr)
       \longrightarrow0. \label{eq:far}
\end{align}
Thus both terms in \eqref{eq:split} tend to zero along the same
selected indices. This contradicts \eqref{eq:integerlower} and
proves Theorem~\ref{thm:general}, hence also
Theorem~\ref{thm:strict}. \qed

\section*{Verification}

An independent mathematical audit (13 November 2026) found no error or
missing hypothesis in the underlying proof. This was not formal verification.

\begin{thebibliography}{9}
\bibitem{BKKKZ}
K.~Barreto, J.~Kang, S.-h.~Kim, V.~Kova\v{c} and S.~Zhang,
\emph{Irrationality of rapidly converging series: A problem of
Erd\H{o}s and Graham},
Bull.\ Lond.\ Math.\ Soc.\ \textbf{58} (2026), no.~8, e70457.
\href{https://doi.org/10.1112/blms.70457}{doi:10.1112/blms.70457}.
\href{https://arxiv.org/abs/2601.21442v3}{arXiv:2601.21442v3}.

\bibitem{Erdos1975}
P.~Erd\H{o}s,
\emph{Some problems and results on the irrationality of the sum of
infinite series},
J.\ Math.\ Sci.\ \textbf{10} (1975), 1--7.
\href{https://www.renyi.hu/~p_erdos/1976-44.pdf}{Full text}.

\bibitem{ErdosGraham}
P.~Erd\H{o}s and R.~L.~Graham,
\emph{Old and new problems and results in combinatorial number theory},
Monographies de L'Enseignement Math\'ematique, vol.~28,
Universit\'e de Gen\`eve, L'Enseignement Math\'ematique, Geneva, 1980.
\href{https://mathweb.ucsd.edu/~ronspubs/80_11_number_theory.pdf}{Full text}.
\end{thebibliography}

\end{document}
